\exercisesection{5}

\begin{enumerate}
\def\labelenumi{\arabic{enumi}.}
\tightlist
\item
  Let A be a commutative ring and \(\mathfrak{I}\) an ideal of A. An A-module M is called absolutely Hausdorff with respect to \(\mathfrak{I}\) if, for every finitely generated A-module N, the A-module N \(\otimes_{\mathrm{A}}\) M is Hausdorff with the \(\mathfrak{I}\) -adic topology; an absolutely Hausdorff A-module is ideally Hausdorff.
\end{enumerate}

\begin{enumerate}
\def\labelenumi{(\alph{enumi})}
\item
  For M to be absolutely Hausdorff with respect to \(\mathfrak{I}\) , it is necessary and sufficient that, for every finitely generated A-module N and every submodule \(\mathbf{N}'\) of N, \(\operatorname{Im}(\mathbf{N}' \otimes_{\mathbf{A}} \mathbf{M})\) be closed in \(\mathbf{N} \otimes_{\mathbf{A}} \mathbf{M}\) with the \(\mathfrak{I}\) -adic topology.
\item
  Let B be a commutative A-algebra and \(\mathfrak{L}\) an ideal of B containing \(\mathfrak{I}B\) . If M is an absolutely Hausdorff B-module with respect to \(\mathfrak{L}\) , M is an absolutely Hausdorff A-module with respect to \(\mathfrak{I}\) .
\end{enumerate}

\begin{enumerate}
\def\labelenumi{\arabic{enumi}.}
\setcounter{enumi}{1}
\item
  Show that every Z-module which is Hausdorff with the p-adic topology (p a prime number) is ideally Hausdorff, but give an example of a finitely generated Z-module which is Hausdorff but not absolutely Hausdorff with respect to p (use Exercise 1 (a)).
\item
  With the notation of Exercise 11 of §3, let N be the submodule of \(\hat{\mathbf{E}}\) which is the closure in \(\hat{\mathbf{E}}\) of the submodule of \(\hat{\mathbf{E}}\) generated by the vectors \(pe_{2n-1} - p^n e_{2n} (n \geqslant 1)\) and let M = \(\hat{\mathbf{E}}/\mathbf{N}\) . Show that under the \(p\) -adic topology the submodule \(p\mathbf{M}\) of M is not closed in M and deduce that M is a \(\mathbf{Z}_p\) -module which is ideally Hausdorff but not absolutely Hausdorff with respect to \(p\) .
\end{enumerate}

¶ 4. Let A be a commutative ring, \(\mathfrak{I}\) an ideal of A and S = 1 + \(\mathfrak{I}\). Show that, if M is an absolutely Hausdorff A-module with respect to \(\mathfrak{I}\), then S \(^{-1}\) M = M, in other words, for all \(a \in\) S, \(x \mapsto ax\) is a bijection of M onto itself. (Show first that the hypothesis that M is Hausdorff with the \(\mathfrak{I}\)-adic topology implies that \(x \mapsto ax\) is injective; then prove that the submodule aM of M is dense in M with the \(\mathfrak{I}\)-adic topology and use Exercise 1 (a).)

\begin{enumerate}
\def\labelenumi{\arabic{enumi}.}
\setcounter{enumi}{4}
\tightlist
\item
  Take \(\mathbf{A} = \mathbf{Z}\) and \(\mathfrak{I} = p\mathbf{Z}\) , where \(p\) is prime.
\end{enumerate}

\begin{enumerate}
\def\labelenumi{(\alph{enumi})}
\item
  Show that, if \(q\) is a prime number distinct from \(p\) , the \(\mathbf{Z}\) -module \(\mathbf{Z} / q\mathbf{Z}\) satisfies property (ii) of Theorem 1, but not property (i).
\item
  Show that the \(\mathbf{Z}\) -module \(\mathbf{Q} / \mathbf{Z}\) satisfies property (iv) of Theorem 1 but not property (ii).
\end{enumerate}

\begin{enumerate}
\def\labelenumi{\arabic{enumi}.}
\setcounter{enumi}{5}
\tightlist
\item
  Let A be a commutative Noetherian ring, \(\mathfrak{I}\) an ideal of A and M an A-module. Show that condition (v) of Theorem 1 is equivalent to the following:
\end{enumerate}

(v') For every finitely generated A-module N and every submodule N' of N, the canonical mapping \((\mathbf{M} \otimes_{\mathbf{A}} \mathbf{N}')^{\wedge} \to (\mathbf{M} \otimes_{\mathbf{A}} \mathbf{N})^{\wedge}\) (where the two sides are the Hausdorff completions with respect to the \(\mathfrak{I}\) -adic topologies on M \(\otimes_{A} N'\) and M \(\otimes_{A} N\) respectively) is injective. (To prove that (v) implies (v'), argue as in part (E) of the proof of Theorem 1. To see that (v') implies (v), consider A-modules N annihilated by a power of \(\mathfrak{I}\).)

¶ 7. Let \((A_{\lambda}, f_{\mu\lambda})\) be a directed direct system of Noetherian local rings; if \(m_{\lambda}\) is the maximal ideal of \(A_{\lambda}\) , suppose that, for \(\lambda \leqslant \mu\) , \(m_{\mu} = m_{\lambda}A_{\mu}\) and that \(A_{\mu}\) is a flat \(A_{\lambda}\) -module. Show then that \(A = \lim A_{\lambda}\) is Noetherian and a flat \(A_{\lambda}\) -module for all \(\lambda\) . (If \(m = m_{\lambda}A\) is the maximal ideal of \(A\) (Chapter II, § 3, Exercise 16), show that on \(A\) the m-adic topology is Hausdorff, observing that \(A\) is a faithfully flat \(A_{\lambda}\) -module for all \(\lambda\) . Then prove that, if \(\hat{A}\) is the completion of \(A\) with respect to the m-adic topology, \(\hat{A}\) is Noetherian, using § 2, no. 10, Corollary 5 to Theorem 2; finally, prove that, for all \(\lambda\) , \(\hat{A}\) is a flat \(A_{\lambda}\) -module using Theorem 1 of no. 2 and Proposition 2 of no. 3.)

¶ 8. Let A be a Noetherian local ring, m its maximal ideal and k = A/m its residue field. Let K be an extension of k; show that there exists a local homomorphism from A to a Noetherian local ring B such that B/mB is isomorphic to K and B is a flat A-module. (Consider first the case where \(\mathbf{K} = k(t)\) , distinguishing two cases according to whether t is algebraic or transcendental over k; consider next a family \((\mathrm{K}_{\lambda})\) of subfields of K containing k, which is well ordered by inclusion and such that, if \(K_{\lambda}\) has a predecessor \(\mathrm{K}_{\mu}, \mathrm{K}_{\lambda} = \mathrm{K}_{\mu}(t_{\mu})\) for some \(t_{\mu} \in K_{\mu}\) . Finally, apply Exercise 7.)

CHAPTER IV(*)

